\documentclass[12pt]{article}
\usepackage[utf8]{inputenc}
\usepackage[T1]{fontenc}
\usepackage[margin=1in]{geometry}
\usepackage{setspace}
\usepackage{xcolor}
\usepackage[colorlinks=true, urlcolor=blue!60!black, linkcolor=blue!60!black]{hyperref}
\setlength{\parindent}{0em}
\setlength{\parskip}{0.6em}
\renewcommand{\baselinestretch}{1.0}
\begin{document}
\begin{center}
{\LARGE \textit{West Virginia v. B. P. J. / Little v. Hecox}, Nos. 24-43 \& 24-38 (2026)\par}
\vspace{0.5em}
{\large Case Brief --- Two-Page Edition\par}
\end{center}

\section{Memory Jogger}
States may limit public-school women's and girls' teams to biological females under both Title IX and the Equal Protection Clause, without individualized exceptions (\href{https://www.law.cornell.edu/supremecourt/text/24-43}{Cornell LII}).

\section{Detailed Case Facts}
West Virginia (2021) and Idaho (2020) barred biological males from female school teams. B. P. J., a transgender girl who took puberty blockers and estrogen and never underwent male puberty, was barred from girls' cross-country and track. Lindsay Hecox, a transgender woman who underwent male puberty, sought to run for Boise State's women's teams.

\section{Procedural History}
The Fourth Circuit ruled for B. P. J. on Title IX and remanded her equal protection claim for factfinding on athletic advantage; the Ninth Circuit affirmed an injunction for Hecox. The Supreme Court reversed both on June 30, 2026.

\section{Judicial Votes}
9--0 in judgment on Title IX; 6--3 on Equal Protection. Kavanaugh wrote for Roberts, Thomas, Alito, Gorsuch, and Barrett; Thomas and Gorsuch concurred. Sotomayor, with Kagan and Jackson, concurred in part and dissented in part; Jackson also wrote separately.

\section{Holding}
``Sex'' in Title IX means biological sex, so schools may keep female teams for biological females. The laws classify by sex, not transgender status, and satisfy intermediate scrutiny because they serve safety and competitive fairness; no individualized exemptions are required. \textit{Hecox} was not moot.

\section{Analysis of Opinions}
\textbf{Majority:} The 1972 meaning of ``sex,'' the Javits Amendment, and the 1975 regulations support separate teams, and sports are ``generally zero sum.'' An as-applied challenge is a ``backdoor argument for strict scrutiny,'' and courts should defer amid ``ongoing medical and scientific debate.'' \textbf{Thomas:} transgender status is not a suspect class. \textbf{Gorsuch:} Title IX gave no clear notice forbidding separate teams, and \textit{Bostock} is consistent. \textbf{Sotomayor:} the Court should have allowed factfinding on a subclass who never underwent male puberty, and it diluted intermediate scrutiny. \textbf{Jackson:} Title IX protects gender identity.

\section{Examples: Future Applications}
\textbf{Same-Side:} a biological-sex rule for girls' swimming is upheld even as applied to a girl who took puberty blockers. \textbf{Opposite-Side:} expelling a student for transitioning remains sex discrimination under \textit{Bostock}. \textbf{Fence-Sitter:} a State law that \textit{includes} transgender girls on girls' teams, left open in footnote 2.

\section{Critique}
Textualists praise fidelity to the 1972 meaning of ``sex'' and democratic federalism. Erwin Chemerinsky warns that the Court watered down intermediate scrutiny for all sex discrimination, and Craig Konnoth argues that plaintiffs litigated the wrong issue.

\section{Key Quotations}
\begin{quote}
``Consistent with Title IX and the Equal Protection Clause, we hold that the States may maintain women's and girls' sports for biological females. They may determine eligibility for women's and girls' sports based on biological sex. The Constitution and Title IX do not require an overhaul of women's and girls' sports throughout America.''
\end{quote}
\noindent --- Justice Kavanaugh, for the Court, slip op. at 29.

\end{document}
